\documentclass[a4paper,12pt]{jarticle}

\usepackage{amsmath,amssymb,amsthm}
\usepackage{mathtools}
\usepackage{bm}
\usepackage{xurl}
\urlstyle{same}

\setlength{\textwidth}{160mm}
\setlength{\textheight}{235mm}
\setlength{\oddsidemargin}{0mm}
\setlength{\evensidemargin}{0mm}

% ---- 定理型環境 ----
\theoremstyle{definition}
\newtheorem{definition}{定義}
\newtheorem{assumption}{仮定}
\newtheorem{notation}{記法}
\theoremstyle{plain}
\newtheorem{prop}{命題}
\newtheorem{cor}{系}
\newtheorem{thm}{定理}
\theoremstyle{definition}
\newtheorem{ex}{例}
\theoremstyle{remark}
\newtheorem*{remark}{注意}

\renewcommand{\proofname}{証明}
\renewcommand{\qedsymbol}{$\blacksquare$}

% ---- 記号 ----
\newcommand{\Act}{S}   % 起動中の能力集合 S(t)
\newcommand{\GF}{\mathrm{GF}}    % ギャラリーフェイク
\newcommand{\OS}{\mathrm{OS}}    % オーダースタンプ
\newcommand{\BV}{\mathrm{BV}}    % ブラックボイス
\newcommand{\CH}{\mathrm{CH}}    % 転校生（コンバートハンズ）
\newcommand{\SM}{\mathrm{SM}}    % サンアンドムーン

\title{クロロ対ヒソカ戦におけるコンボの理解\footnote{Ver1.0（2026年7月12日）：本稿は論文の形式を借りたファン考察・パロディの一種であり，原作の描写との整合性を保つよう心がけていますが，あくまで洒落なのでガチで突っ込むのはご遠慮ください．とはいえ，明白な矛盾についてはできる限り修正するつもりですので，追加すべき命題と共に，ご指摘いただけると幸いです．
なお富樫先生は34巻末尾の解説で「ストーリーを作る際のマニュアルは自分の中にあって，それは全て言語化できますし，実際考えた末，マニュアル通りが面白いという結論に達することも多々あります」と述べています．「マニュアル」がどこまでを指すのか不明ですが，富樫先生の頭の中には本稿の命題群が無矛盾に整理されている可能性が高いです．まことにもって驚くべき知性というほかないでしょう．こういった優れた作品を母国語で読めることに感謝します．}\\
\large ― 形式的体系化の試み ―}
\author{浜田　宏\\東北大学大学院\footnote{文学研究科行動科学専攻分野教授}}
\date{}

\begin{document}
\maketitle

\begin{abstract}
 『HUNTER × HUNTER（ハンターハンター）』34 巻におけるクロロ対ヒソカ戦では，クロロが盗んだ複数の能力による連携コンボが展開され，物語の大きな見せ場となっている． 
対ヒソカ戦におけるクロロ＝ルシルフルの戦術は，「同時に起動できるのは2能力」という制約のもとで，
複数の盗んだ能力を死後残る念を利用して連携させる点に独創性がある．
本稿では，各能力を定義し，その効果をフロー型／ストック型に分類したうえで各種仮定を明示し，
戦闘中の各コンボを命題として定式化し，それぞれに証明を与える．
複雑な戦術を理解する鍵は起動数と効果の持続の区別にあり，フロー型の産物（複製）をストック型の刻印（死後の念）で固定して次フェーズへ持ち越すことで，起動枠を空けながらコンボが成立する点にある．
これにより，クロロが3つ以上の能力を同時に起動することなくヒソカ戦に勝利したこと，
および一見同時使用に見える場面（ブラックボイスとサンアンドムーンの併用など）が矛盾しないことを示す．
さらに考察として，こうしたコンボ設計が，作者・冨樫義博の
カードゲームへの深い造詣（グリードアイランド編に顕著）を土台にしている可能性についても論じる．
\end{abstract}

\bigskip


\section*{はじめに}
『HUNTER $\times$ HUNTER（ハンターハンター）』34巻におけるクロロ対ヒソカ戦では，クロロの能力「盗賊の極意（スキルハンター）」を用いた華麗な戦闘コンボが展開され，物語の大きな見せ場の一つとなっている．

しかし多くの読者と同様に，本稿の筆者も，何度読み返しても理解できず，「まあなんかよくわからんけど，天才富樫先生のことだから，多分つじつまはあってるんだろう」くらいに考え，完全な理解を諦めていた．あるとき，この難解さは，数学のテキストをナナメ読みしていても一向に頭に入らない現象と同型であることにふと気づいた．筆者の職業は数理社会学者であり，普段数学を含むテキストを読むときには必ず手元に紙とペンを置き，定義・定理・証明を追いながらメモを取ることを習慣としている．そのようにして初めて，数学の命題はその内容が理解できるのである．

つまり『HUNTER $\times$ HUNTER（ハンターハンター）』には，メモを取りながら読まないと完全に理解できないレベルの複雑な命題が含まれていることに気づいたのである（特に物語の後半）．

ならばいっそ数学の命題の証明をフォローするように，クロロの戦術を形式的に整理してみよう，普段やっている研究と同じ水準の厳密さで『HUNTER $\times$ HUNTER（ハンターハンター）』に向き合おう，と考えたのが本稿を執筆するに至った動機である．

本稿を読んでいる読者の皆さんには，最新刊が出る度に王位継承戦の頭から読み直し，その度に34巻で止まってしまったという経験はないだろうか．理由の一端は，クロロの戦術の複雑さにある．少なくともこの戦闘に関する疑問を解消し，心置きなく復習してから最新刊（本稿執筆時点では39巻）を楽しむための一助となれば幸いである．

%======================================================================
\section{定義：各能力}

\begin{notation}[起動中の能力集合]
時刻 $t$ において起動中の能力（スキル）の集合を $\Act(t)$ と書く．
能力 $s$ が時刻 $t$ で使用可能であることを $s\in\Act(t)$ と表す．
\end{notation}

\begin{notation}[効果]
能力 $s$ が生み出す物理的な結果（複製・刻印・命令・変装など）を，単に $s$ の効果とよぶ．
\end{notation}

\begin{definition}[盗賊の極意（スキルハンター）]
クロロの本体能力（特質系）．盗んだ念能力を本の各ページに封じ込める．
封じた能力 $s$ は，$s\in\Act(t)$ のときに限り使用できる．

基本状態（栞を用いない場合）では，封じた能力を使うには
本を開き，右手で持つ必要がある（定義\ref{def:active}(i)）．
\end{definition}

スキルハンターの使用中は本を開いて右手で持つため片手がふさがり，両手を要する能力の使用が困難になる．
栞（ダブルフェイス，定義\ref{def:active}(ii)）は，本を持たなくてもそのページの能力を使えるようにする，拡張手段である．

\begin{definition}[起動と栞のテーマ（ダブルフェイス）]\label{def:active}
能力$s$ が起動中（$s\in\Act(t)$）であるとは，次のいずれかが成り立つことをいう：
\begin{enumerate}
\item[(i)] クロロが本を右手に持ち，能力$s$ のページを開いている；
\item[(ii)] 能力$s$ のページに栞が挟まれている（ダブルフェイス）．
\end{enumerate}
ここで，本を開くことで使用できる能力は基本的には一つであり，栞は1枚である．なお，栞は本のページに挟まれている間だけ有効であり，栞を外すと挟んでいた頁の能力は停止する．
\end{definition}

\begin{definition}[効果の型]\label{def:type}
能力 $s$ の効果を次のいずれかに分類する：
\begin{itemize}
\item フロー型：$s\in\Act(t)$ である間だけ効果が持続し，$s$ を停止した瞬間に消滅する．
\item ストック型：一度発生した効果は，以後 $s\notin\Act(t)$ となっても持続する．
\end{itemize}
\end{definition}

\begin{definition}[人形]\label{def:doll}
オーダースタンプ（定義\ref{def:os}）による操作の対象となりうる存在を{人形}とよぶ．
オーダースタンプの元の持ち主が「人形である」と認識できるもの，すなわち胴体と首が繋がった人型がこれにあたる．
胴体と首が分離した瞬間に，その対象は人形でなくなる．
\end{definition}


なお，死体そのものは人形ではない．
死体は胴体と首が繋がった人型ではあるが，オーダースタンプの操作対象にならない（定義\ref{def:os}(i)）．
死体をギャラリーフェイクで複製したもの（死体のコピー）が，はじめて人形となる．

\begin{quotation}
オレからすれば死体も人形も同じただの動かない塊なのだが，オーダースタンプの元の持ち主だった者にすれば違うらしい．$\cdots$前の持ち主の心の闇に思いを馳ながら能力を自分のものにしてゆく．スキルハンターの醍醐味だよ（30-32頁）．
\end{quotation}

また，生きた人間も人形ではなく，その操作にはブラックボイスを用いる（定義\ref{def:bv}）．
まとめると，
\[
\text{生者}\ (\text{人形でない})，\quad
\text{死体}\ (\text{人形でない})，\quad
\text{死体のコピー}\ (\text{人形})．
\]


\begin{definition}[神の左手悪魔の右手（ギャラリーフェイク，$\GF$）]\label{def:gf}
左手で触れた物体の複製を右手側に生成する．生物・念製物は複製できない
（複製できるのは物体，または死体に限る）．
この複製（人形）が GF の効果である．
\end{definition}

なお『神の左手悪魔の右手』は楳図かずお先生の作品タイトルに，ルビのギャラリーフェイクは細野不二彦先生の作品タイトルに由来すると考えられる．原作中でも屈指のかっこよさを誇る命名である．
なおこの能力名に読みと無関係のルビを振る手法は，荒木飛呂彦先生（ジョジョの奇妙な冒険）もしばしば使っている\footnote{余談ではあるが漫画界における超能力描写の三大発明は大友克洋『童夢』，荒木飛呂彦『ジョジョの奇妙な冒険』，富樫義博『ハンターハンター』だと個人的に思っている．この三作品は後に大量のエピゴーネンを生み出したという意味でも記念的偉業である．}．


\begin{definition}[番いの破壊者（サンアンドムーン，$\SM$）]\label{def:sm}
左手で太陽（$+$）の刻印を，右手で月（$-$）の刻印を対象に押す．
相異なる符号の刻印が接触すると爆発する．
\end{definition}



\begin{definition}[人間の証明（オーダースタンプ，$\OS$）]\label{def:os}
スタンプを押した人形（定義\ref{def:doll}）に命令を与えて操作する．次の制約をもつ：
\begin{enumerate}
\item[(i)] 死体そのものは操作できない．ただし死体の複製（ギャラリーフェイク，定義\ref{def:gf}）は操作できる．
\item[(ii)] 人形の胴体と首が分離すると，その対象は人形でなくなり，操作は解け，押されたスタンプも消える（定義\ref{def:doll}）．
\item[(iii)] 人間をコピーした人形に「殺せ」と命じると原型の性格を反映して非効率化・命令無視（フリーズ）が生じうるが，「壊せ」と命じると一律に相手の首を胴体から外そうとする．
\end{enumerate}
\end{definition}

『人間の証明』は森村誠一の長編推理小説もしくは松田優作主演映画として映像化された同小説の角川映画タイトルに由来すると考えられる．


\begin{definition}[携帯する他人の運命（ブラックボイス，$\BV$）]\label{def:bv}
アンテナを刺した生きた人間の意思を捻じ曲げ，命令通りに操作する．
\end{definition}

\begin{definition}[転校生（コンバートハンズ，$\CH$）]\label{def:ch}
右手で触れた相手をクロロの姿に，左手で触れた対象の姿にクロロが変わる（変装）．
\end{definition}

大林宣彦監督作『転校生』のタイトルに由来すると考えられる．広島・尾道を舞台に身体が入れ替わってしまった男女の中学生を描いた青春映画である．なお新海誠監督作『君の名は』は言うまでもなく，この映画のオマージュである．

\begin{cor}[姿の交換]\label{cor:swap}
クロロが同一の対象 $X$ を両手で触れると，クロロと $X$ の姿は入れ替わる．
\end{cor}

\begin{proof}
定義\ref{def:ch} より，右手が触れた $X$ はクロロの姿になり，
左手が触れた $X$ の姿にクロロが変わる．
両手で同一の $X$ に触れれば，$X$ はクロロの（元の）姿を，
クロロは $X$ の姿をとるから，両者の姿は交換される．
\end{proof}

\begin{definition}[能力集合]\label{def:P}
以上をまとめ，クロロがヒソカ戦で用いた能力全体の集合を
\[
\mathcal{S}=\{\text{スキルハンター},\ \text{ダブルフェイス},\ \GF,\ \SM,\ \OS,\ \BV,\ \CH\}
\]
と書く．スキルハンター（盗賊の極意）は本体能力であり，ダブルフェイスは本に付属する能力と説明されている．
その他は本に封じた各ページ（盗んだ能力）に対応する．
起動集合はこの部分集合，すなわち $\Act(t)\subseteq\mathcal{S}$ を満たす．
\end{definition}

なお，この集合の要素数は，クロロが一人の相手に対して用いた能力数の最大値である．

\begin{quotation}
 一人殺るに当たり使用する能力数新記録だ．誇りに思っていいよ（25頁）． 
\end{quotation}

%======================================================================
\section{仮定}

\begin{assumption}[ギャラリーフェイクはフロー型]\label{as:gf}
GF による複製はフロー型である（定義\ref{def:type}）．
すなわち $\GF\notin\Act(t)$ となった瞬間に複製は消滅する．
\end{assumption}

\begin{assumption}[サンアンドムーンはストック型：死後強まる念]\label{as:sm}
SM の刻印はストック型であり，一度押すと爆発するまで消えない．
（元の持ち主が死亡した「死後強まる念」であることによる．）
\end{assumption}

\begin{assumption}[刻印の威力と接触時間]\label{as:power}
サンアンドムーンの爆破の威力（爆破範囲）は接触時間に依存する．
原作のクロロの説明によれば，威力を最大化するには対象（人間）に $3$〜$5$ 秒触れ続ける必要がある．
これを次のように敷衍する：一瞬の接触なら爆破範囲は刻印の周辺のみ，
触れ続けるほど範囲は徐々に広がり，最終的に人体全体に及ぶ．
必要な時間は人体の大きさに比例する（小柄なら約 $3$ 秒，大柄なら約 $5$ 秒で全体）．
また刻印の対象となりうるのは人間（者）に限ると仮定する．
（この解釈のもとで，審判の腹部が爆発したのは背中側の刻印の爆破が貫通した結果と説明できる．
原作76--80頁に対応する仮定．）
\end{assumption}

\begin{assumption}[ブラックボイスはフロー型]\label{as:bv}
BV による（生きた人間の）操作はフロー型である．
すなわち操作の継続には $\BV\in\Act(t)$（栞を挟むか本を開く，定義\ref{def:active}）が必要であり，
$\BV\notin\Act(t)$ となると操作は解ける．
なお1つの $\BV$ の起動で複数の人間を同時に操作でき（原作44頁：2人を同時操作），
操作に用いるアンテナは念ではなく実体である．
サンアンドムーンと異なり，ブラックボイスの持ち主（シャルナーク）は本戦闘の時点で存命であり
「死後強まる念」ではない．したがって栞を外す・本を閉じると能力は持ち主へ返り，
効果は持続しないと解する
（ダブルフェイスの制約とブラックボイスの仕様について参考文献[\ref{ref:huncyclo}][\ref{ref:kirua}]）．
\end{assumption}

\begin{assumption}[転校生（変装）はフロー型]\label{as:ch}
CH による変装（偽クロロの姿）はフロー型である．姿を維持するには
$\CH\in\Act(t)$（栞を挟むか本を開く）が必要であり，$\CH\notin\Act(t)$ となると元の姿に戻る．
（クロロが潜伏中に栞を転校生に挟み続ける描写に基づく．）
\end{assumption}

\begin{assumption}[刻印による固定]\label{as:anchor}
ストック型の刻印を付与された対象は，その存在自体がストック化する．
特に，$\GF$ の複製に $\SM$ の刻印を押すと，$\GF\notin\Act(t)$ となっても複製は消滅しない．
（原作88頁：コピーと刻印は「死後の念」として残る，とヒソカが到達する．）
\end{assumption}

\begin{assumption}[オーダースタンプはフロー型]\label{as:agent}
OS による（人形の）操作はフロー型である．人形の制御を継続するには
$\OS\in\Act(t)$（栞を挟むか本を開く，定義\ref{def:active}）が必要であり，
$\OS\notin\Act(t)$ となると操作は解ける．
（原作29頁に，本の頁をめくるとオーダースタンプの効果が消える描写があり，
フロー型であることの直接の傍証となる．）
ただし1つの $\OS$ の起動で，スタンプを押したすべての人形を同時に制御できる
（人形が何体でも起動枠は1つ）．
\end{assumption}

\begin{remark}
仮定\ref{as:anchor}・\ref{as:bv}・\ref{as:power} は原作の明示が薄く，考察により解釈が分かれる
暫定的な仮定である（一方，仮定\ref{as:agent} は原作29頁の描写に裏づけられる）．
以下では上述の定義・仮定のもとで内部整合的な体系を構成する．
\end{remark}

%======================================================================
\section{命題：戦闘中のコンボ}

以下，戦闘中に用いられた各コンボを命題として定式化し，証明する．
各命題には，対応する原作の場面を例として付す
（頁数はコミックス34巻）．

%----------------------------------------------------------------------
\begin{prop}[同時起動の上限（定義\ref{def:active} の帰結）]\label{prop:two}
定義\ref{def:active} から，任意の時刻 $t$ について $|\Act(t)|\le 2$ が従う．
すなわちクロロは高々2つの能力しか同時に起動できず，
3つ以上を同時に起動することはできない．
\end{prop}

\begin{proof}
定義\ref{def:active} より，能力が起動中であるのは
「本を開いている」か「栞が挟まれている」場合に限る．
本は同時に1ページしか開けないので (i) による起動は高々1つ，
栞は1枚であるから (ii) による起動も高々1つ．
ゆえに $|\Act(t)|\le 2$．
この上限は独立の仮定ではなく，定義\ref{def:active} から導かれる帰結である．
\end{proof}

\begin{remark}[起動数と効果の持続は別物]
命題\ref{prop:two} の $|\Act(t)|\le2$ は，各時刻に起動中
（定義\ref{def:active}：本を開くか栞を挟む）の能力数の上限であって，
すでに発生した効果の個数の上限ではない．
ストック型の効果（サンアンドムーンの刻印，およびそれに固定された複製；
仮定\ref{as:sm}・\ref{as:anchor}）は，能力を停止したのちも持続するが，
これらは $\Act(t)$ に数えられない．
したがって，多数の刻印人形が盤上に残っていても，
同時に{起動}している能力が2つ以下であれば命題\ref{prop:two} と矛盾しない．
この区別（起動数 $\neq$ 効果の持続）が，以降のコンボ整合性の議論全体を支える．
\end{remark}

%----------------------------------------------------------------------
\begin{cor}[両手動作の制約]\label{cor:twohands}
あるスキルの発動動作が両手を要するなら，その動作の瞬間において，
\begin{enumerate}
\item[(1)] 当該スキルは栞で固定されていなければならず（本を持てないため），
\item[(2)] 本は閉じていなければならない（両手を空けるため），
\item[(3)] したがって起動中のスキルはそれ1つだけである：$\Act(t)=\{s\}$，$|\Act(t)|=1$．
\end{enumerate}
\end{cor}

\begin{proof}
定義\ref{def:active} より，スキル $s$ が起動中であるのは
(i) 本を開いて手に持つ，または (ii) 栞を挟む場合に限る．
両手動作中は本を持てないので (i) は成立せず，$s$ は (ii) 栞で固定されていなければならない．
かつ両手を空けるため本を閉じる．本を閉じている間は栞のスキルのみが起動中であり，
栞は1枚（定義\ref{def:active}）だから，起動中のスキルは $s$ ただ1つである．
\end{proof}

\begin{remark}
「両手を要する」はスキルではなく動作ごとに判定する．
たとえばサンアンドムーンは，$+$ 刻印だけを押すなら片手で足りる
（本を開いたまま，他スキルと併用可）が，$+$ と $-$ を同時に配置するなら両手を要する．
ギャラリーフェイクの複製作成（左手で触れ右手から出す）や，
転校生の姿の交換（系\ref{cor:swap}）は両手動作である．
この区別により，命題\ref{prop:mass} のステップ（$\Act=\{\GF,\SM\}$）は本系と矛盾しない——
そこではサンアンドムーンを片手で押しているだけで，ギャラリーフェイクは維持のみだからである．
なお本系はクロロが単独で本を扱う場合の結論であり，
本を外部エージェントに持たせる運用は考慮しない．
\end{remark}

%----------------------------------------------------------------------
\begin{prop}[複製の消滅と刻印による固定]\label{prop:anchor}
ギャラリーフェイクの複製は，栞（ダブルフェイス）で固定していない限り，
本を閉じれば消える．
ただしサンアンドムーンの{刻印}を押すことで，本を閉じても消えない．
\end{prop}

\begin{proof}
まず前半を示す．仮定\ref{as:gf} より GF による複製はフロー型である．
複製を栞で固定していない場合，本を閉じると定義\ref{def:active} の (i)(ii) の
いずれも満たされず $\GF\notin\Act(t)$ となる．
フロー型の効果は $\GF\notin\Act(t)$ の瞬間に消滅する（定義\ref{def:type}）．
ゆえに複製は消える．

次に後半を示す．複製に $\SM$ の刻印を押したとする．
仮定\ref{as:sm} より刻印はストック型で，爆発するまで消えない．
仮定\ref{as:anchor} より，刻印を付与された複製はストック化する．
したがって $\GF\notin\Act(t)$（本を閉じた状態）でも複製は消滅しない．
\end{proof}

\begin{ex}
原作において，クロロは審判の死体をギャラリーフェイクで複製したうえで，
その複製にサンアンドムーンの刻印を押している．
これにより，戦闘後半を通じて多数の人形がギャラリーフェイクの起動なしに存在し続ける．
「刻印を押された複製が消えずに残る」というこの描写は，命題\ref{prop:anchor} の後半の一例である（原作30頁・88頁）．
\end{ex}

\begin{remark}[原作におけるミスリード（33頁・88頁）]
原作33頁でクロロはヒソカに，「人形（複製）が残るのはギャラリーフェイクを解除していないためで，
人形が消えればギャラリーフェイクを解除したと思ってよい」と開示する．
これは複製の存続をギャラリーフェイクの\emph{維持}に帰する説明だが，ミスリードである．
88頁でヒソカが到達する真の仕組みは本命題のとおりで，
コピーと刻印は「死後の念」として残る（ゆえにギャラリーフェイクは解除済みでよい）．
クロロは「複製が残る $\Leftrightarrow$ GF維持中」と誤認させることで，
実際には空いている $\GF$ の起動枠を隠し，敵の推理を撹乱している．
\end{remark}

%----------------------------------------------------------------------
\begin{prop}[刻印人形の量産と解放]\label{prop:mass}
クロロは，任意の個数の刻印付き複製人形を生成でき，
かつ生成後はギャラリーフェイクを停止（解放）してよい．
この一連の操作は各時刻 $|\Act(t)|\le 2$ を保って実行できる．
\end{prop}

\begin{proof}
次の手順をとる．
\begin{enumerate}
\item $\GF$ に栞を挟んで起動する（$\Act=\{\GF\}$，$|\Act|=1$）．複製を作成する．
\item $\SM$ を本で開いて起動する（$\Act=\{\GF,\SM\}$，$|\Act|=2$）．
      各複製に刻印を押す．
\item $\GF$ と $\SM$ を停止する（$\Act=\varnothing$）．
\end{enumerate}
各ステップで $|\Act(t)|\le 2$．
命題\ref{prop:anchor} より，刻印を押された複製は手順3の後も残る．
複製生成は繰り返せるので，任意個の刻印人形を得たうえで $\GF$ を解放できる．
\end{proof}

\begin{ex}
原作で，クロロは審判の死体をギャラリーフェイクでコピーし，のちに観客の複製人形を大量に製造する．
これは，複製生成を繰り返して多数の人形を得るという命題\ref{prop:mass} の量産部分の一例である．
\end{ex}

%----------------------------------------------------------------------
\begin{prop}[コピー人形と材料との同一性]\label{prop:copy}
観客のコピー人形は，生きた観客からではなく，
サンアンドムーンの爆発で生じた死体をギャラリーフェイクで複製して作られる．
同じ顔の人形が多数並ぶのは同一の死体を複製したためであり，
転校生（コンバートハンズ）によるものではない．
\end{prop}

\begin{proof}
定義\ref{def:gf} より，ギャラリーフェイクは生物を複製できず，
複製できるのは物体または死体に限る．
ゆえにコピー人形の材料は死体でなければならない．
死体はサンアンドムーンの爆発（開戦時の審判の爆殺など）によって生じる．
また，ギャラリーフェイクの複製は元の対象と同一の外見をもつから，
同一の死体を種とすれば，生成される人形はすべて同じ顔になる．
一方，転校生は1つの対象の姿を別の姿へ変換する能力であって，
体を増やす（量産する）能力ではない．
ゆえに同一顔の人形が多数存在することは転校生では説明されず，
ギャラリーフェイクによる同一死体の複製として説明される．
\end{proof}

\begin{ex}
原作で，観客席に同じ顔の人物が多数見られるのは，本命題の帰結の一例である．
クロロは戦闘で生じた審判・観客の死体をギャラリーフェイクで複製した（原作34巻30頁）．
\end{ex}

% \begin{remark}[数と準備に関する議論：共闘説]
% 本命題はコピーの{作り方}を説明するが，その{数と準備の速さ}については議論がある．
% ギャラリーフェイクが生物を複製できないこと（定義\ref{def:gf}），
% ダブルフェイスにおいて栞を外すと能力が持ち主（コルトピ）へ返ること，
% 複製が栞または刻印なしには消えること（命題\ref{prop:anchor}）を踏まえると，
% 観客席に見られる多数（200体規模とも言われる）の同顔人形を
% クロロ単独で用意できたかは自明でない．
% このことから，コルトピ（ギャラリーフェイクの持ち主で，無制限に量産できる）や
% シャルナーク（多人数の操作が可能）が関与したとする{共闘説}が提起されている
% （参考文献[\ref{ref:fakutsu}]）．
% 単独説か共闘説かは，作者コメントを含めても決着していない．
% \end{remark}

%----------------------------------------------------------------------
\begin{prop}[人形化の連鎖]\label{prop:chain}
生きた人間から，オーダースタンプで操作可能な人形を直接得ることはできない．
人形を得るには，
\[
\text{対象}\ \xrightarrow{\ \SM\ \text{等で死体化}\ }\ \text{死体}
\ \xrightarrow{\ \GF\ \text{で複製}\ }\ \text{死体のコピー}
\ \xrightarrow{\ \OS\ \text{でスタンプ}\ }\ \text{操作可能な人形}
\]
という3段階を要する．
\end{prop}

\begin{proof}
定義\ref{def:gf} より $\GF$ は生物を複製できないから，複製の前に対象を死体化する必要がある（第1段階）．
定義\ref{def:os}(i) より $\OS$ は死体そのものを操作できないが，死体の複製は操作できる．
ゆえに死体を直接スタンプするのではなく，一度 $\GF$ で複製してから（第2段階）
その複製にスタンプを押す（第3段階）必要がある．
なお，生きた人間は人形（定義\ref{def:doll}）ではないため $\OS$ の対象外であり，
生者の操作にはブラックボイスを用いる（命題\ref{prop:agent}）．
\end{proof}

%----------------------------------------------------------------------
\begin{prop}[首の分離による人形の停止]\label{prop:decap}
人形の胴体と首が分離すると，その対象は人形でなくなり，
オーダースタンプの操作は解け，押されていたスタンプも消える．
ただし{複製の体そのものは死後の念により残る}（消えない）．
\end{prop}

\begin{proof}
定義\ref{def:doll} より，胴体と首の分離をもって対象は人形でなくなる．
定義\ref{def:os} より $\OS$ の操作対象は人形に限られるから，
人形でなくなった対象は操作対象から外れ，スタンプも消える．
一方，複製の存続は刻印（仮定\ref{as:anchor}）が支えるので，首を外しても体は残る
（{原作33頁}：首を外された人形はスタンプが消えるが複製は残る）．
\end{proof}

\begin{ex}
原作で，クロロが人形へ与える命令が「殺せ」ではなく\textbf{「壊せ」}であるのは，
本命題の制約に由来する（命題\ref{prop:fight} の例を参照）．
定義\ref{def:os}(iii) の通り「殺せ」では原型の性格が出て非効率化・フリーズが生じうるため，
クロロは「壊せ」と命じた．「壊せ」の命令下で人形は一律に相手（ヒソカ）の首を胴体から外そうとする．
逆にヒソカの側も，人形の首を落とせばその人形を停止させられる（命題\ref{prop:decap}）．
（原作34巻32頁）
\end{ex}

%----------------------------------------------------------------------
\begin{prop}[操作と変装の両立（3能力同時は不要）]\label{prop:fight}
命題\ref{prop:mass} で固定された人形を用いれば，
オーダースタンプによる人形操作と，転校生による自身の変装とを，
$|\Act(t)|\le 2$ を保って同時に行える．
このとき{ギャラリーフェイクは起動していない}．
\end{prop}

\begin{proof}
命題\ref{prop:mass} より，人形の存在維持に $\GF$ は不要であり $\GF\notin\Act(t)$．
よって戦闘フェーズで起動しうるのは高々 $\{\OS,\CH\}$ の2つで，$|\Act(t)|\le2$．
ゆえに「$\GF$，$\OS$，$\CH$ の3能力を同時起動する」必要はない．
（人形攻撃中に $\GF$ が動いているように見えるのは，
命題\ref{prop:anchor} により複製が刻印で固定され，$\GF$ から{切り離されている}ためである．）
\end{proof}

\begin{ex}
原作で，クロロは量産した複製人形にオーダースタンプを押し，
「ヒソカを壊せ」と命令して一斉に攻撃させる．
このとき複製はすでに刻印で固定されており（命題\ref{prop:anchor}），
ギャラリーフェイクは起動していない．
人形を操作しつつ自身は群衆に紛れるこの場面は，命題\ref{prop:fight} の一例である（原作34巻92頁）．

\end{ex}

%----------------------------------------------------------------------
\begin{prop}[外部エージェントによる手の拡張]\label{prop:agent}
クロロは，操作下のエージェントに物理タスクを肩代わりさせ，
能力起動に必要な「手」や「動作主」を実質的に増やせる．
オーダースタンプ・ブラックボイスはいずれもフロー型で，操作の継続に起動枠を1つ要する
（仮定\ref{as:agent}・\ref{as:bv}）．
ただしオーダースタンプは，1つの起動枠で{スタンプを押した全人形を同時に}制御できるため，
1枠で多数の「手」を動かせる．
\end{prop}

\begin{proof}
仮定\ref{as:agent}・\ref{as:bv} より，オーダースタンプ・ブラックボイスはフロー型であり，
操作の継続には起動（$\OS\in\Act(t)$，$\BV\in\Act(t)$）が必要で，起動枠を1つ消費する．
ただし物理タスク自体（攻撃・保持・移動）は操作対象が実行するため，クロロ自身の手は空く．
特にオーダースタンプは1枠で全人形を同時制御できるので，1枠から多数の動作主が得られる．
\end{proof}

\begin{ex}
原作の開戦と同時に，クロロは審判（生きた人間）にブラックボイスのアンテナを刺して操作し，
ヒソカへ攻撃を差し向ける．これはブラックボイスで生者を操作する例であり，
その継続にはブラックボイスの起動（栞または本の開口）を要する（仮定\ref{as:bv}）（原作19頁・44頁）．
\end{ex}

%----------------------------------------------------------------------
\begin{prop}[自爆トリガーの構成]\label{prop:bomb}
太陽（$+$）刻印と月（$-$）刻印を別々の対象に配置し，
オーダースタンプで人形を動かして両刻印を接触させることで，
クロロは任意の時点・任意の位置で爆発を起こせる．
\end{prop}

\begin{proof}
定義\ref{def:sm}（サンアンドムーン）より，相異なる符号の刻印が接触すると爆発する．
命題\ref{prop:anchor} より，刻印は $\SM$ 停止後も対象に残り続ける．
命題\ref{prop:fight} より，$\OS$ 単独（$|\Act|\le2$）で人形を移動させられる．
したがって，$+$ 刻印を持つ人形と $-$ 刻印を持つ対象を接触させる命令を与えれば，
その瞬間に爆発が生じる．発火タイミングはクロロが命令によって制御できる．
\end{proof}

\begin{ex}
原作で，操作された審判の腹部が突然爆発する場面がサンアンドムーンの初発動であり，
あらかじめ仕込まれた刻印による起爆の一例である（参考文献[\ref{ref:warachie}]）．
また終盤，ヒソカが手にしていた頭部が突然爆発する決着も，
刻印を仕込んだ対象を接触・起爆させたものであり，命題\ref{prop:bomb} の決定的な適用例である（原作34巻19頁，105頁）．
\end{ex}

\begin{cor}[人形と人間の同時操作による起爆]\label{cor:dualbomb}
$+$ 刻印を持つ人形をオーダースタンプで，$-$ 刻印を持つ生きた人間をブラックボイスで
別々に操作し，両者を接触させれば起爆できる．
このとき起動は $\{\OS,\BV\}$ の2つで $|\Act(t)|\le2$．
\end{cor}

\begin{proof}
仮定\ref{as:agent}・\ref{as:bv} より，オーダースタンプで人形を，ブラックボイスで生きた人間を，
それぞれ起動枠1つずつ（計2つ）で操作できる．
$+$ 刻印の人形と $-$ 刻印の人間を接触させれば，命題\ref{prop:bomb} により起爆する．
\end{proof}

\begin{ex}
原作104--105頁の決着で，クロロは月刻印（$-$）を押した人間をブラックボイスで操作し，
{ヒソカの手に頭部が戻った瞬間}を見計らって，太陽刻印（$+$）のついた胴体に接触させて爆破する
（参考文献[\ref{ref:kazunoko}]）．これは本系の一例である．
\end{ex}

\begin{remark}[ブラックボイスとサンアンドムーンの「見かけの同時使用」]
原作19頁で，ブラックボイスで操作した審判の腹部がサンアンドムーンで爆破される．
一見この2能力を同時に使っているように見えるが，そうではない．
サンアンドムーンはストック型（仮定\ref{as:sm}）であり，
{刻印を「押す」にはSMの起動が要るが，「爆発」には起動を要さない}
（$+$ と $-$ の接触で自動的に起爆する）．すなわち爆発は{効果の持続}であって
{起動}ではない（命題\ref{prop:two} 直後の注意）．
したがってこの瞬間に起動しているのはBVだけであり，
刻印は{あらかじめ}（SMを起動していた別の時点で）押されている．
サンアンドムーンは起動していないので「両手が要る・本を閉じる」制約はこの瞬間には生じず，
栞をSMに移す必要もない．ゆえに矛盾しない．
この{事前刻印}による理解はファン考察とも一致する（参考文献[\ref{ref:warachie}][\ref{ref:fakutsu}]）．
\end{remark}

%----------------------------------------------------------------------
\begin{prop}[本体の秘匿]\label{prop:hide}
転校生による変装と，多数の操作人形への紛れ込みにより，
ヒソカはどれが本体のクロロかを一意に特定できない．
\end{prop}

\begin{proof}
定義\ref{def:ch}（転校生）より，クロロは自身の姿を任意の人物の姿に変えられる．
命題\ref{prop:mass} より，同一の場に多数の人形が存在する．
変装したクロロと人形は外見上区別できないから，
観測者（ヒソカ）は本体の位置を確率的にしか推定できない．
このとき $\CH$ の起動は1枠のみを消費し，他の1枠を戦術に回せる（命題\ref{prop:two}）．
\end{proof}

\begin{ex}
原作で，クロロは転校生によって他人（観客）の姿になり，群衆に紛れる．
戦闘はヒソカの視点からのみ描かれ，読者もヒソカも本体のクロロを特定できない．
これは命題\ref{prop:hide} の一例である．

さらに作中には，コンバートハンズ（定義\ref{def:ch}）で{クロロの姿に変えられた観客}が登場し，
ヒソカの攻撃を受けて（本体でないことが）{正体を暴かれる}場面がある（原作84頁・85頁）．
\end{ex}

%----------------------------------------------------------------------
\begin{prop}[偽クロロの製造と操作]\label{prop:fake}
クロロは，転校生で観客を{自分（クロロ）の姿}に変えて{偽クロロ}を作り，
ブラックボイスでそれを操作してヒソカと交戦させられる．
この運用は $|\Act(t)|\le2$ を保って成立する．
\end{prop}

\begin{proof}
転校生（定義\ref{def:ch}）の右手側の効果により，触れた観客をクロロの姿に変えて偽クロロを得る．
このとき左手で自分自身も変装するなら両手を要するので，
生成の瞬間は本を閉じ転校生のみが起動する（系\ref{cor:twohands}）．
生成後，偽クロロの姿の維持には転校生の起動が必要（仮定\ref{as:ch}）なので，転校生に栞を挟む．
次にブラックボイスのページを開き，偽クロロにアンテナを刺して操作する（仮定\ref{as:bv}）．
このとき起動は $\{\CH,\BV\}$ の2つで $|\Act(t)|\le2$．
偽クロロがヒソカと戦う間，本体のクロロは別に潜伏でき（命題\ref{prop:hide}），
ヒソカは偽と本体を区別できない．
\end{proof}

\begin{ex}
原作で，クロロは観客を偽クロロに変えてブラックボイスで操作し，ヒソカと交戦させる．
これはコピー人形（物量）とは別系統の{囮}のコンボであり，本命題の一例である
（参考文献[\ref{ref:kazunoko}]）．
また原作45頁からは，栞をCHへ移すとブラックボイスが解ける（仮定\ref{as:bv}）ため，
クロロは{BVで操作した人間が用済みになってから}栞をCHへ移した，と読み取れる（順序の制約, 原作34巻84頁）．
\end{ex}

%======================================================================
\section{定理：コンボ全体の整合性}

\begin{thm}[運用全体の無矛盾性]\label{thm:main}
「複製を作り $\to$ 刻印で固定し $\to$ 人形に命令して攻撃させ $\to$ 自身は変装し $\to$
任意時点で爆発を起こす」という一連のコンボは，
命題\ref{prop:two}（同時2つまで）に反することなく実行できる．
\end{thm}

\begin{proof}
運用を{準備フェーズ}（命題\ref{prop:mass}）と
{戦闘フェーズ}（命題\ref{prop:fight}・\ref{prop:bomb}）に時間分割する．
命題\ref{prop:mass} より準備フェーズは $|\Act(t)|\le2$，
命題\ref{prop:fight}・\ref{prop:bomb} より戦闘フェーズも $|\Act(t)|\le2$ で実行できる．
フェーズ間で必要な「複製・刻印の持ち越し」は命題\ref{prop:anchor} の永続性が保証する．
また命題\ref{prop:agent} により，物理動作は操作対象へ委譲できる．
ゆえに全時刻で $|\Act(t)|\le2$ を保ったままコンボが成立する．
\end{proof}

\begin{cor}
「ギャラリーフェイク $+$ オーダースタンプ $+$ 転校生」を同時に3つ起動する必要はない．
戦闘フェーズでは $\GF\notin\Act(t)$ だからである（命題\ref{prop:fight}）．
\end{cor}

\begin{cor}
オーダースタンプ・転校生・ブラックボイスがいずれもフロー型（常時起動が必要）であっても，
コンボは成立する．戦闘フェーズで実際に起動するのは $\{\OS,\CH\}$ のちょうど2つであり，
命題\ref{prop:two} に反しない．
ここで $\GF$ を手放せるのは，サンアンドムーンの刻印が複製を固定して
{存在の枠（$\GF$）を解放している}ためである（命題\ref{prop:anchor}・\ref{prop:mass}）．
すなわち，人形が消えないのはオーダースタンプではなく{刻印}の効果であり，
これによって $\GF$ の枠が空く．
\end{cor}

\begin{cor}[コンボの本質]
本コンボの核心は
\[
\underbrace{\text{GFの複製（フロー型）}}_{\text{放置すると消える}}
\ \xrightarrow{\ \SM\ \text{の刻印（ストック型）で固定}\ }\
\underbrace{\text{永続する刻印人形}}_{\text{フェーズをまたいで持ち越せる}}
\]
という{効果の型変換}にある．
すなわち「消える能力の産物を，消えない能力で固定して次フェーズへ持ち越す」点に，
同時2つの制約を回避する鍵がある．
\end{cor}

%======================================================================
\section{考察}

本稿のように，クロロの戦術を定義・仮定・命題・証明の形式へ落とし込んでみると，
その設定が驚くほど{内部整合的}であることに改めて気づかされる．
各能力は「発動条件」「効果の持続（フロー型／ストック型）」「対象の制約」を備え，
それらの組合せ（コンボ）が矛盾なく機能する．
これは，設定が場当たり的に継ぎ足されたものではなく，
一貫したルール体系のうえに設計されていることを強く示唆する．

この論理性の源泉として，作者である冨樫義博先生の
{カードゲームへの深い造詣}を挙げることができる．
実際，グリードアイランド編では，カードの種別・ランク・複製・保管・発動条件などを
精緻に規定した，非常によく練られたカードシステムが作中に登場する．
興味深いことに，本稿で用いた語彙——「効果の持続」「複製」「固定（永続化）」
「発動枠の上限」——は，そのままカードゲームのルール用語に読み替えられる．
たとえば，フロー型／ストック型の区別は使い捨てカードと永続（パーマネント）カードの区別に，
刻印による複製の固定は永続効果やカウンターの付与に，
同時起動数 $|\Act(t)|\le2$ は場・手札・コストの制約に，それぞれ対応づけられる．

以上を踏まえると，クロロの能力コンボは，
{カードゲームにおけるカードスキルのコンボ}のアナロジーによって非常にすっきりと理解できる．

盗賊の極意（スキルハンター）が，他者の能力を「カード」のように
収集・保管し，場に出して連携させる構造そのものが，
デッキ構築型カードゲームの発想と地続きなのである．
本稿が示した「コンボが同時2つの制約下で無矛盾に成立する」という事実は，
その背後に，カードゲーム的に整備されたルール設計が存在することの
一つの傍証といえるだろう．

%======================================================================
\section*{注意}

\noindent
本稿の骨格（定義群と仮定\ref{as:gf}--\ref{as:sm}）は原作描写とファン考察に整合するが，
仮定\ref{as:anchor}（刻印が複製を固定する），
仮定\ref{as:bv}（ブラックボイスの継続条件と返却）および
仮定\ref{as:power}（刻印の威力と接触時間）は原作での明示が薄く，
複数の解釈が成立する余地がある．したがって仮定の一部には原作者の意図通りかどうか不明な箇所が存在する（それゆえに「仮定」している）．
本稿は「原作から確定した定理」ではなく，
「妥当な仮定のもとで内部整合的な体系」を構築したパロディ論文として読んでいただきたい．

\bigskip
\section*{参考文献}
{\small

本稿を作成する上で参考にした資料・考察記事を以下に列挙すると共に，熱い情報を公共財として公開してくれたことを（同じ原作ファンとして）感謝する．また，原作の描写はコミックス34巻（2017年6月30日発行第一刷）を参照した．



\begin{enumerate}
\item\label{ref:kazunoko} かずのこ「【クロロvsヒソカ】何が起きていた？戦いの流れを徹底解説！」note．\\
  \url{https://note.com/djkazunoko/n/n75265488fef8}（2026年7月12日取得）．
\item\label{ref:hedorah} 「351話から356話まででクロロが何をしたか」ハンター$\times$ハンターの考察とか．\\
  \url{https://hedorah123.hatenablog.com/entry/2016/06/15/134408}（2026年7月12日取得）．
\item\label{ref:fakutsu} ふぁくつ「HUNTER$\times$HUNTER考察 ヒソカvsクロロ 4つの問いに答える」note．\\
  \url{https://note.com/ks_kaz_game/n/n63b2ef29dd63}（2026年7月12日取得）．

\item\label{ref:huncyclo} 「クロロの念能力『栞のテーマ（ダブルフェイス）』の制約と誓約について考察」HUNTER$\times$HUNTER考察．\\
  \url{https://huncyclopedia.com/335}（2026年7月12日取得）．
\item\label{ref:mocho} 「ヒソカVSクロロ戦について～能力解説編」もちょ分が！．\\
  \url{https://mou-chotto.com/2017/07/07/manga-hunter-quwrof_vs_hyskoa1/}（2026年7月12日取得）．
\item\label{ref:intwosick} 「クロロVSヒソカ！ハンターハンター34巻前半のあらすじ・ネタバレ」マンガマスタリー．\\
  \url{https://intwosick.com/hunter34/}（2026年7月12日取得）．
\item\label{ref:asufazuru} アスファズル「人間の証明（オーダースタンプ）」アスファズルの庵のブログ．\\
  \url{https://ameblo.jp/asufazuru/entry-12508112248.html}（2026年7月12日取得）．
\item\label{ref:warachie} 「《ハンターハンター》ヒソカ対クロロ 究極のデスマッチを解説してみた」warachie．\\
  \url{https://warachie.net/hunter-death-match/}（2026年7月12日取得）．
\item\label{ref:kirua} 「携帯する他人の運命（ブラックボイス）」キルアから見たHUNTER$\times$HUNTER．\\
  \url{https://kirua-hunterhunter.com/blackvoice/}（2026年7月12日取得）．
\end{enumerate}
}

\end{document}
